% 第五章 结果讨论与研究局限

\section{结果讨论与研究局限}

\subsection{结果讨论}

第一阶段结果表明，宏观变量可以用于构造条件预期收益率基准，但样本外预测增益整体有限。5日窗口的样本外改进幅度仅为$R^2_{OS}=0.0019$，60日窗口的$R^2_{OS}=0.0392$，虽优于历史均值基准，但提升仍不大。由此看，模型在解释宏观信息如何进入收益率方面，比在样本外预测方面更有价值。第二阶段显示，短期偏离更容易被市场状态变量解释，iVIX与日内振幅在5日窗口中表现稳定，说明风险感知与交易拥挤程度会明显影响短期偏离。长窗口结果对模型设定更敏感，资金与流动性变量在嵌套回归中能提供增量解释力，LASSO则倾向于保留更稀疏的变量集合。扩展模型IV的结果进一步表明，政策不确定性与汇率波动在60日窗口中的解释力增量显著（调整后$R^2$从0.0306提升至0.1165），长期偏离并不只是由市场内部交易状态驱动，也受到政策预期与外部环境变化的共同影响。

本文对“预期偏离”的界定只有一个基准，即第一阶段MIDAS模型生成的条件预期收益率。第二阶段中的iVIX并不参与偏离定义，而是作为风险感知变量，用于解释市场偏离这一基准时为何会被放大。

\subsection{研究局限}

本文仍有几处限制需要明确写出。宏观变量之间存在较强联合线性依赖，多变量MIDAS设定因此难以形成稳定主模型，第一阶段最终收束为相对简化的条件预期收益率构造。第一阶段样本外预测增益整体有限。5日窗口改进幅度很弱，60日窗口虽出现正的样本外改进，但增益仍然有限，故相关结果更适合用于预期基准构造，而不宜扩展为强预测结论。第二阶段在长期窗口下对模型筛选方式较为敏感，部分变量的作用稳定性还需要进一步检验。MIDAS、分组嵌套回归与LASSO都是成熟方法，文章的价值主要落在研究框架组织与变量整合上。文中没有单独展开完整的稳健性结果展示，后续研究可围绕权重函数替换、滞后长度调整、非重叠收益窗口与子样本重估进一步补充稳健性检验，本文正文不再展开这些扩展结果。

扩展模型IV表明，政策不确定性与汇率波动已被正式纳入偏离识别框架，并在60日窗口中带来明显的解释力提升。需要注意的是，EPU单独系数在两个窗口下都未表现出稳定显著性，汇率波动仅在60日窗口显著，故更稳妥的结论应落在新增变量组的联合效应，而不宜将任何单一变量表述为长期偏离的绝对主导因素。